%%%%%%%%%%%%%%%%%%%%%%%%%%%%%%%%%%%%%%%%%%%%%%%%%%%%%%%%%%%%%%%%%%%%%%
% Lớp 10 | Phương pháp tọa độ trong mặt phẳng (𝑂𝑥𝑦) | Ba đường conic trong mặt phẳng tọa độ
% Dạng 0H9?5-2: Phương trình chính tắc của elip
% Số câu: 19  (sinh tự động bởi tach_cau.py; từ nay thầy sửa trực tiếp file này)
%%%%%%%%%%%%%%%%%%%%%%%%%%%%%%%%%%%%%%%%%%%%%%%%%%%%%%%%%%%%%%%%%%%%%%

%%%%%%%%%%%%%%%% TRẮC NGHIỆM 4 PHƯƠNG ÁN %%%%%%%%%%%%%%%%

% nguồn: MaHoaNH_TracNghiem_Lop10
\begin{ex}%[0H9N5-2]%[ma:9dee7041]
	Cho elip $(E)$ có độ dài trục lớn  bằng $8$, độ
	dài trục bé bằng $6$. Phương trình của $(E)$  là
	phương trình nào sau? 
	\choice{$\dfrac{x^2}{64}+\dfrac{y^2}{36}=1$}{$\dfrac{x^2}{8}+\dfrac{y^2}{6}=1$}{\True  $\dfrac{x^2}{16}+\dfrac{y^2}{9}=1$}{$\dfrac{x^2}{16}-\dfrac{y^2}{9}=1$}
	\loigiai{ Ta có  $\heva{&2a=8\\& 2b=6}\Leftrightarrow \heva{&a=4\\&b=3.}$\\
		Vậy phương  trình chính tắc của elip là $\dfrac{x^2}{16}+\dfrac{y^2}{9}=1$.
	} 
\end{ex}

% nguồn: MaHoaNH_TracNghiem_Lop10
\begin{ex}%[0H9N5-2]%[ma:e0f13aa9]
	Cho elip $(E)$ có độ dài trục lớn là $20$, độ dài trục bé là $12$. Lập phương trình chính tắc  của $(E)$.
	 \choice
		{  $\dfrac{x^2}{100}+\dfrac{y^2}{64}=1$}
		{$\dfrac{x^2}{64}+\dfrac{y^2}{100}=1$} 
		 {\True $\dfrac{x^2}{100}+\dfrac{y^2}{36}=1$}
		 {$\dfrac{x^2}{400}+\dfrac{y^2}{144}=1$}
	\loigiai{ Ta có $\heva{&2a=20\\& 2b=12}\Leftrightarrow  \heva{&a=10\\&b=6.}$\\
		Do đó do đó phương trình chính tắc của $(E)$  là $\dfrac{x^2}{100}+\dfrac{y^2}{36}=1$.
	 }
\end{ex}

% nguồn: MaHoaNH_TracNghiem_Lop10
\begin{ex}%[0H9N5-2]%[ma:3fe3e8b1]
	Phương trình nào  sau đây là phương trình chính tắc của một đường e-líp?
	\choice 
		{$\dfrac{x^2}{16}-\dfrac{y^2}{25}=1$} 
		 {\True $\dfrac{x^2}{16}+\dfrac{y^2}{9}=1$}
		 {$\dfrac{x^2}{25}-\dfrac{y^2}{9}=-1$}
		{$\dfrac{x^2}{25}+\dfrac{y^2}{9}=-1$} 
	 \loigiai{
		Phương  trình chính tắc của e-líp là $\dfrac{x^2}{16}+\dfrac{y^2}{9}=1$.
	} 
\end{ex}

% nguồn: MaHoaNH_TracNghiem_Lop10
\begin{ex}%[0H9N5-2]%[ma:57b2b538]
	Phương trình nào sau đây không phải là phương trình  chính tắc của một elip?
	\choice 
		{$\dfrac{x^2}{36}+\dfrac{y^2}{9}=1$} 
		 {$\dfrac{x^2}{5}+\dfrac{y^2}{2}=1$}
		 {$\dfrac{x^2}{25}+\dfrac{y^2}{16}=1$}
		{\True $\dfrac{x^2}{9}-\dfrac{y^2}{4}=1$} 
	\loigiai{ 
		Phương trình elip có dạng $\dfrac{x^2}{a^2}+\dfrac{y^2}{b^2}=1$ với $a>b>0$  nên phương trình $\dfrac{x^2}{9}-\dfrac{y^2}{4}=1$ không phải là phương trình chính tắc của elip.
	} 
\end{ex}

% nguồn: MaHoaNH_TracNghiem_Lop10
\begin{ex}%[0H9N5-2]%[ma:ebe39788]
	Phương trình chính tắc của elip có  độ dài trục lớn bằng 10, độ dài trục bé bằng 8 là
	\choice 
		{\True  $\dfrac{x^2}{25}+\dfrac{y^2}{16}=1$}
		{$\dfrac{x^2}{25}+\dfrac{y^2}{9}=1$} 
		 {$\dfrac{x^2}{100}+\dfrac{y^2}{64}=1$}
		 {$\dfrac{x^2}{16}+\dfrac{y^2}{25}=1$}
	 \loigiai{
	Vìđộ dài trục lớn bằng  $10$ và độ dài trục bé bằng $8$ nên $\heva{&2a=10\\&2b=8}\Leftrightarrow\heva{&a=5\\&b=4.}$\\
	Phương  trình chính tắc của elip là $\dfrac{x^2}{25}+\dfrac{y^2}{16}=1$.
	} 
\end{ex}

% nguồn: MaHoaNH_TracNghiem_Lop10
\begin{ex}%[0H9N5-2]%[ma:a4a2778a]
	Cho elip $(E)$ đi qua  $2$ điểm $A_1(-3;0)$, $B_1(0;-2)$. Phương trình nào là phương trình chính tắc của $(E)$?
	 \choice
		 {$\dfrac{x^2}{4}+\dfrac{y^2}{9}=1$}
		{\True $\dfrac{x^2}{9}+\dfrac{y^2}{4}=1$} 
		{$\dfrac{x^2}{9}-\dfrac{y^2}{4}=1$} 
		 {$\dfrac{x^2}{4}-\dfrac{y^2}{9}=1$}
	\loigiai{ 
		Giả sử phương trình $(E)\colon \dfrac{x^2}{a^2}+\dfrac{y^2}{b^2}=1$, với $a,\,  b>0$. \\
		Vì $(E)$ đi qua hai điểm $A_1(-3;0)$, $B_1(0;-2)$ nên ta có $a^2=9$, $b^2=4$.  \\
		Vậy  $(E)\colon \dfrac{x^2}{9}+\dfrac{y^2}{4}=1$.
	 }
\end{ex}

% nguồn: MaHoaNH_TracNghiem_Lop10
\begin{ex}%[0H9H5-2]%[ma:c4986e31]
	Trong mặt phẳng $Oxy$, phương trình chính tắc của đường  elip với $a=4$, $c=\sqrt{7}$ là
	 \choice
		 {$\dfrac{x^2}{16}+\dfrac{y^2}{9}=0$}
		 {\True $\dfrac{x^2}{16}+\dfrac{y^2}{9}=1$}
		{$\dfrac{x^2}{9}+\dfrac{y^2}{16}=1$} 
		{$\dfrac{x^2}{16}-\dfrac{y^2}{9}=1$} 
	\loigiai{Ta có $b^2=a^2-c^2=16-7=9$, suy ra phương trình chính  tắc của đường elip là $\dfrac{x^2}{16}+\dfrac{y^2}{9}=1$.}
\end{ex}

% nguồn: MaHoaNH_TracNghiem_Lop10
\begin{ex}%[0H9H5-2]%[ma:df41f130]
	Cho elip $(E)$ có một đỉnh $A_2(7;0)$, một tiêu điểm $F_2(5;0)$. Lập phương trình  chính tắc của $(E)$.
	 \choice
		{ \True $\dfrac{x^2}{49}+\dfrac{y^2}{24}=1$} 
		 {$\dfrac{x^2}{24}+\dfrac{y^2}{49}=1$}
		{  $\dfrac{x^2}{49}+\dfrac{y^2}{25}=1$}
		{$\dfrac{x^2}{49}-\dfrac{y^2}{24}=1$} 
	\loigiai{Ta có đỉnh $A_2(7;0)$, một  tiêu điểm $F_2(5;0)$, suy ra $a=7$, $c=5$ và $b^2=a^2-c^2=24$, do đó $\dfrac{x^2}{49}+\dfrac{y^2}{24}=1$.
	 }
\end{ex}

% nguồn: MaHoaNH_TracNghiem_Lop10
\begin{ex}%[0H9H5-2]%[ma:1f2fb778]
	Cho elip $(E)$ có  một đỉnh $A_1(-4;0)$, một tiêu điểm $F_2(3;0)$. Lập phương trình chính tắc của $(E)$.
	 \choice
		{ \True $\dfrac{x^2}{16}+\dfrac{y^2}{7}=1$} 
		 {$\dfrac{x^2}{16}+\dfrac{y^2}{9}=1$}
		{  $\dfrac{x^2}{9}+\dfrac{y^2}{7}=1$}
		{$\dfrac{x^2}{16}-\dfrac{y^2}{5}=1$} 
	\loigiai{Ta có đỉnh $A_1(-4;0)$, một tiêu điểm $F_2(3;0)$, suy ra $a=4$, $c=3$ và $b^2=a^2-c^2=7$, do đó  $\dfrac{x^2}{16}+\dfrac{y^2}{7}=1$.
	 }
\end{ex}

% nguồn: MaHoaNH_TracNghiem_Lop10
\begin{ex}%[0H9H5-2]%[ma:fee46c1d]
	Cho elip $(E)$ có độ dài trục lớn là $6$, độ dài  tiêu cự là $4$. Lập phương trình chính tắc của $(E)$.
	\choice 
		 { \True $\dfrac{x^2}{9}+\dfrac{y^2}{5}=1$}
		 {$\dfrac{x^2}{36}+\dfrac{y^2}{16}=1$}
		{  $\dfrac{x^2}{5}+\dfrac{y^2}{9}=1$}
		 {$\dfrac{x^2}{9}-\dfrac{y^2}{5}=1$}
	\loigiai{ 
		Ta có $\heva{&2a=6\\& 2c=4}\Leftrightarrow \heva{&a=3\\&c=2}\Rightarrow b^2=a^2-c^2=5$, do đó do đó phương trình chính tắc của $(E)$ là  $\dfrac{x^2}{9}+\dfrac{y^2}{5}=1$.
	} 
\end{ex}

% nguồn: MaHoaNH_TracNghiem_Lop10
\begin{ex}%[0H9H5-2]%[ma:3cd23cec]
	Cho  elip $(E)$ có độ dài trục bé là $6$, độ dài tiêu cự là $4$. Lập phương trình chính tắc của $(E)$.
	 \choice
		 { \True $\dfrac{x^2}{13}+\dfrac{y^2}{9}=1$}
		{$\dfrac{x^2}{36}+\dfrac{y^2}{16}=1$} 
		{ $\dfrac{x^2}{9}+\dfrac{y^2}{4}=1$} 
		 {$\dfrac{x^2}{9}+\dfrac{y^2}{13}=1$}
	\loigiai{Ta có $\heva{&2b=6\\& 2c=4}\Leftrightarrow  \heva{&b=3\\&c=2}\Rightarrow a^2=b^2+c^2=13$.\\
		Do đó phương trình chính tắc  của $(E)$ là $\dfrac{x^2}{13}+\dfrac{y^2}{9}=1$.
	} 
\end{ex}

% nguồn: MaHoaNH_TracNghiem_Lop10
\begin{ex}%[0H9H5-2]%[ma:11d6a238]
	Cho elip $(E)$ đi qua một điểm $M\left(-3\sqrt{2};\dfrac{3\sqrt{2}}{2}\right)$, một tiêu điểm  $F_1\left(-3\sqrt{3};0\right)$. Lập phương trình chính tắc của $(E)$.
	 \choice
		 { \True $\dfrac{x^2}{36}+\dfrac{y^2}{9}=1$}
		 {$\dfrac{x^2}{25}+\dfrac{y^2}{9}=1$}
		{ $\dfrac{x^2}{36}+\dfrac{y^2}{25}=1$} 
		{$\dfrac{x^2}{36}-\dfrac{y^2}{9}=1$} 
	\loigiai{Ta có một tiêu  điểm $F_1(-3\sqrt{3};0)$, suy ra $c=3\sqrt{3}$ và $a^2-b^2=27$.\\
		Vì $M \in (E)$  nên $\dfrac{18}{a^2}+\dfrac{9}{2b^2}=1$.\\
		Ta có hệ phương trình $\heva{&a^2-b^2=27\\&\dfrac{18}{a^2}+\dfrac{9}{2b^2}=1} \Leftrightarrow \heva{&a^2=36  \\&b^2=9.}$\\
		Do đó phương trình chính tắc  của $(E)$ là $\dfrac{x^2}{36}+\dfrac{y^2}{9}=1$.
	 }
\end{ex}

% nguồn: MaHoaNH_TracNghiem_Lop10
\begin{ex}%[0H9H5-2]%[ma:cc5d0eef]
	Trong mặt phẳng $Oxy$, phương trình chính tắc của Elip có độ dài trục lớn bằng $6$ và tỉ số giữa tiêu cự với độ dài trục lớn bằng $\dfrac{1}{3}$  là
	 \choice
		{$\dfrac{x^2}{9}+\dfrac{y^2}{3}=1$} 
		{$\dfrac{x^2}{9}+\dfrac{y^2}{5}=1$} 
		{\True  $\dfrac{x^2}{9}+\dfrac{y^2}{8}=1$}
		{$\dfrac{x^2}{6}+\dfrac{y^2}{5}=1$} 
	\loigiai{ Độ dài trục lớn $2a=6\Rightarrow a=3$.\\ 
		Tỉ số giữa tiêu cự và độ  dài trục lớn $\dfrac{2c}{2a}=\dfrac{1}{3}\Rightarrow c =\dfrac{1}{3} \cdot a=\dfrac{1}{3}\cdot 3=1$.\\
		Ta có $c^2=a^2-b^2\Leftrightarrow 1^2=3^2-b^2\Leftrightarrow b^2=3^2-1^2=8$. \\ 
		Phương trình chính tắc của Elip $\dfrac{x^2}{9}+\dfrac{y^2}{8}=1$. 
	} 
\end{ex}

% nguồn: MaHoaNH_TracNghiem_Lop10
\begin{ex}%[0H9H5-2]%[ma:4c7dd2d5]
	Trong mặt phẳng $Oxy$,  phương trình chính tắc của elip có độ dài trục lớn là $34$ và tiêu điểm $F_2(8;0)$ là
	 \choice
		{$\dfrac{x^2}{17^2}+\dfrac{y^2}{8^2}=1$} 
		{$\dfrac{x^2}{34^2}+\dfrac{y^2}{30^2}=1$} 
		 {$\dfrac{x^2}{15^2}+\dfrac{y^2}{8^2}=1$}
		{\True $\dfrac{x^2}{17^2}+\dfrac{y^2}{15^2}=1$} 
	 \loigiai{
		\begin{itemize}
			\item Độ dài  trục lớn $2a=34\Leftrightarrow a=17$.
			\item Vì tiêu điểm $F_2(8;0)$  nên $c=8$.
			\item Ta có  $b=\sqrt{a^2-c^2}=\sqrt{17^2-8^2}=15$.
		\end{itemize}
		Vậy phương trình elip  là $\dfrac{x^2}{17^2}+\dfrac{y^2}{15^2}=1$.
	} 
\end{ex}

% nguồn: MaHoaNH_TracNghiem_Lop10
\begin{ex}%[0H9V5-2]%[ma:31b6862c]
	Cho elip $(E)$ đi qua một điểm $M\left(5\sqrt{2};4\sqrt{2}\right)$,  một tiêu điểm $F_1(-6;0)$. Lập phương trình chính tắc của $(E)$.
	\choice 
		{  \True $\dfrac{x^2}{100}+\dfrac{y^2}{64}=1$}
		 {$\dfrac{x^2}{64}+\dfrac{y^2}{100}=1$}
		{ $\dfrac{x^2}{100}+\dfrac{y^2}{36}=1$} 
		{$\dfrac{x^2}{400}-\dfrac{y^2}{144}=1$} 
	\loigiai{Ta có một tiêu điểm $F_1=(-6;0)$, suy ra $c=6$ và $a^2-b^2=36$.\\ 
		Vì $M \in (E)$ nên $\dfrac{50}{a^2}+\dfrac{32}{b^2}=1$.  Ta có hệ phương trình $\heva{&a^2-b^2=36\\&\dfrac{x^2}{a^2}+\dfrac{y^2}{b^2}=1} \Leftrightarrow \heva{&a^2=100 \\&b^2=64.}$\\
		 Do đó phương trình chính tắc của $(E)$ là $\dfrac{x^2}{100}+\dfrac{y^2}{64}=1$.
	} 
\end{ex}


%%%%%%%%%%%%%%%% ĐÚNG/SAI %%%%%%%%%%%%%%%%

% nguồn: MaHoaNH_DungSai_Lop10
\begin{ex}%[0H9H5-2]%[ma:5d872fc9]
	Cho elip $(E)$ có tiêu cự bằng $12$ và tổng khoảng cách từ mỗi điểm thuộc elip đó tới hai tiêu điểm của elip bằng $20$. Trong các mệnh đề sau, mệnh đề nào đúng, mệnh đề nào  sai?
	\choiceTF 
		{Các tiêu điểm là $F_1(-12;0)$;  $F_2(12;0)$}
		{\True  Độ dài trục lớn $20$ }
		{\True Phương trình chính tắc của  elip là $\dfrac{x^2}{100}+\dfrac{y^2}{64}=1$}
		{Các giao điểm của $(E)$ với  trục tung là $B_1(-8;0)$; $B_2(8;0)$}
	\loigiai{Gọi phương trình chính tắc của  elip là $\dfrac{x^2}{a^2}+\dfrac{y^2}{b^2}=1$ với $a>b>0$.\\
		Ta có tiêu cự $2c=12$, suy ra  $c=6$ và $2a=20$ nên $a=10$.\\
		Suy ra $b^2=a^2-c^2=100-36=64$. Vậy phương trình chính tắc của elip là  $\dfrac{x^2}{100}+\dfrac{y^2}{64}=1$.
	 }
\end{ex}

% nguồn: MaHoaNH_DungSai_Lop10
\begin{ex}%[0H9H5-2]%[ma:bcfc4c40]
	Cho elip $(E)$ đi qua điểm $A(5;0)$ và có một tiêu điểm là $F_2(3;0)$. Trong các mệnh  đề sau, mệnh đề nào đúng, mệnh đề nào sai?
	 \choiceTF
		{\True Tiêu  điểm $F_1(-3;0)$}
		{\True Tổng các khoảng cách từ mỗi điểm trên elip tới  hai tiêu điểm bằng $10$}
		{\True Tiêu cự của  elip là $F_1F_2=6$}
		{\True Elip có phương  trình chính tắc là $\dfrac{x^2}{25}+\dfrac{y^2}{16}=1$}
	\loigiai{Gọi phương trình chính tắc của  elip là $\dfrac{x^2}{a^2}+\dfrac{y^2}{b^2}=1$ với $a>b>0$.\\
		Ta có $A(5;0)\in (E)$ nên  $a^2=25$.\\
		Lại có $c=3$, suy ra $b^2=a^2-c^2=16$.  Tiêu cự là $2c=6$.\\
		Vậy phương  trình chính tắc của elip là $\dfrac{x^2}{25}+\dfrac{y^2}{16}=1$.
	 }
\end{ex}

% nguồn: MaHoaNH_DungSai_Lop10
\begin{ex}%[0H9H5-2]%[ma:3a3df544]
	Cho elip $(E)$ có phương trình chính tắc là $\dfrac{x^2}{a^2} + \dfrac{y^2}{b^2} = 1$. Biết rằng $(E)$ đi qua điểm $M(\sqrt{5};1)$ và khoảng cách từ một đỉnh nằm trên trục  lớn đến một đỉnh nằm trên trục nhỏ bằng tiêu cự. Trong các mệnh đề sau, mệnh đề nào đúng, mệnh đề nào sai?
	 \choiceTF
		{Điểm $Q(1;1)$  thuộc elip}
		{\True Các tiêu điểm là  $F_1\left(-\dfrac{2\sqrt{6}}{3};0\right)$, $F_2\left(\dfrac{2\sqrt{6}}{3};0\right)$}
		{Tiêu cự  bằng $\dfrac{5\sqrt{6}}{3}$}
		{Phương trình chính tắc của elip là $\dfrac{x^2}{\dfrac{4}{3}}+\dfrac{y^2}{4}=1$} 
	 \loigiai{
		 \noindent
		Gọi  $A(a;0)$, $B(0;b)$ là hai đỉnh.
		Khi đó $AB=2c$, suy ra  $a^2+b^2=4c^2$.\\
		Mà $c^2=a^2-b^2$ nên $3a^2=5b^2$.\\ 
		Vì $M\in (E)$ suy ra $b^2=4$,  $a^2=\dfrac{20}{3} \Rightarrow c=\dfrac{2\sqrt{6}}{3}$.\\
		Vậy  $F_1\left(-\dfrac{2\sqrt{6}}{3};0\right)$, $F_1\left(-\dfrac{2\sqrt{6}}{3};0\right)$.
	} 
\end{ex}


%%%%%%%%%%%%%%%% TRẢ LỜI NGẮN %%%%%%%%%%%%%%%%

% nguồn: MaHoaNH_TraLoiNgan_Lop10
\begin{bt}%[0H9H5-2]%[ma:8b201ad8]
	Trong mặt phẳng $Oxy$, cho Elip $(E)$ có dạng chính tắc. Biết $(E)$ đi qua điểm $A(6;4)$ và có độ dài trục lớn gấp hai lần độ dài trục nhỏ. Tính tiêu cự của $(E)$ (làm tròn kết quả  đến hàng phần mười).
	\shortans{17,3} 
	\loigiai{ 
		Gọi phương trình chính tắc của elip  là $\dfrac{x^2}{a^2}+\dfrac{y^2}{b^2}=1,\,(a>b)$.\\
		Elip đi  qua điểm $A(6;4)$, ta có $\dfrac{36}{a^2}+\dfrac{16}{b^2}=1$.\\
		Elip có độ dài trục lớn gấp hai lần độ dài trục nhỏ, ta có $a=2b$.\\ 
		Ta có hệ phương trình 
		 $$\heva{&\dfrac{36}{a^2}+\dfrac{16}{b^2}=1\\&a=2b}\Leftrightarrow\heva{&a^2=100\\&b^2=25.}$$
		Suy ra  $c^2=a^2-b^2=100-25=75\Rightarrow b=5\sqrt{3}$.\\
		Vậy $2c=2\cdot 5\sqrt{3}=10\sqrt{3}\sim eq 17,3$. 
	} 
\end{bt}
